\chapter{Related Work}
\label{ch:related_work}

\section{The Foundation: Classical Early Warning Signals}
A system heading toward a tipping point changes the way its data looks, and it starts doing so before the tipping point is reached. If those changes can be measured, they act as an early warning. Two papers turned this into a usable method. \citet{scheffer2009} set out the underlying principle and showed it applies across many fields. \citet{dakos2012} fixed it into a step-by-step procedure and released it as software. Nearly every data-based early-warning study since, including \citet{bury2021} which this thesis builds on most directly, either uses that procedure or reacts to it.

\subsection{The Generic Principle: Scheffer et al. (2009)}
\citet{scheffer2009} set out a single principle, drawn from evidence across many fields: whenever a system approaches a bifurcation, it becomes slower to recover from small disturbances. This is critical slowing down, and Chapter~\ref{ch:theoretical_background}, Section~\ref{sec:csd} works through why it happens. The effect is the same in every system, so it leaves the same fingerprint in the data, and that fingerprint can be watched for without knowing the system's equations.

The paper shows this across cases that have nothing else in common: an epileptic seizure beginning, a financial market swinging toward crisis, and a shift from a warm to a cold climate 34 million years ago. In each, the paper documents a rise in variance, in lag-1 autocorrelation, or in both, as the tipping point nears -- the fingerprint predicted by critical slowing down.

Both rise because recovery has slowed. Each random disturbance now carries the system further from its normal level before it returns, which widens the spread of values -- the variance. And each value is measured before the system has fully returned, so it stays close to the previous value -- the autocorrelation. Chapter~\ref{ch:theoretical_background} gives the formulae; what matters here is that neither measurement depends on the specific system, which is what makes a data-only warning possible.

These indicators also have clear limits, and the limits shape this thesis. They can show that a bifurcation may be approaching, but not which kind. They can be set off by other changes: a growing amount of random disturbance, for instance, raises the variance on its own. And they need a record long enough, and finely enough sampled, to measure.

\subsection{The Standard Method: Dakos et al. (2012)}
\label{sec:classical_method}
\citet{dakos2012} turned this principle into a fixed routine, released as an R package called \texttt{earlywarnings}, so later studies could run the same analysis without re-choosing its parts. The routine has four steps.

\begin{enumerate}
  \item \textbf{Detrend.} Estimate the slow-moving background level and subtract it, leaving only the short-term fluctuation around it. The default is Gaussian kernel smoothing, with linear detrending and first-order differencing as alternatives. How much to smooth (the bandwidth) is chosen by the user.
  \item \textbf{Slide a window.} Pick a window length -- half the record is the usual choice -- and move it along the leftover series one step at a time.
  \item \textbf{Measure.} At each window position, compute the target statistic: usually the variance or the lag-1 autocorrelation.
  \item \textbf{Test for a trend.} Use Kendall's $\tau$, a rank correlation between that statistic and time, to summarise whether it rises across the record. $\tau$ runs from $-1$ to $+1$; a value near $+1$ means it rose almost step by step.
\end{enumerate}

Two details from this routine recur throughout this thesis.

The first is the use of Kendall's $\tau$ rather than an ordinary correlation. An ordinary correlation measures how close a rise is to a straight line, but critical slowing down produces a rise that starts slow and steepens. Kendall's $\tau$ only asks whether the values keep increasing, not whether they increase evenly, which is what is needed here.

The second is that $\tau$ alone proves nothing: a short, noisy record can reach a high $\tau$ by chance. \citet{dakos2012} handle this by building many artificial records with the same noise as the real one but no approaching bifurcation, then checking whether the real $\tau$ stands out from them. This thesis uses the same test, in the form \citet{bury2021} later adopted; Chapter~\ref{ch:theoretical_background}, Section~\ref{sec:ar1_surrogates} gives the procedure.

\subsection{Application to This Thesis's Data: Hennekam et al. (2020)}
\label{sec:hennekam_ews}
\citet{hennekam2020} applied the standard sliding-window method (Section~\ref{sec:classical_method}) to the three cores this thesis uses: 64PE406E1 and MS66, both from below 1{,}600~m water depth, and MS21, from about 1{,}000~m. They detrended each molybdenum record with a Gaussian kernel, used a rolling window half the length of the record, and tested the resulting Kendall's $\tau$ trends against 1{,}000 surrogate series.

The result depended on water depth. In the two deep cores, variance rose before every anoxic event analysed, and lag-1 autocorrelation rose before most of them; both trends were far stronger than the surrogates could produce by chance (combined $p < 10^{-13}$ for variance, $p < 0.003$ for autocorrelation). In the shallow core MS21, the signals were largely absent. The same pattern held for a second proxy, uranium, and across the analysis rising variance was the more consistent of the two indicators.

This matters to the thesis in two ways. First, the signal is consistent with a sapropel onset being a real bifurcation, but does not establish one: a gradual change in the driving conditions, with no tipping point, can raise variance and autocorrelation in the same way, and these indicators cannot separate the two cases. Second, they do not identify the type of bifurcation, a gap Section~\ref{sec:classical_limitations} returns to. Re-testing the bifurcation reading against a matched null, and then attempting to name the type on these same cores, is what this thesis's method is built to do.

\citet{bury2021} also used these cores, from the same \citet{hennekam2020} dataset, as one of their real-world tests, and restricted the analysis to molybdenum and uranium -- the two elements that track the oxygen state of the bottom water most directly. This thesis makes the same restriction for its type analysis and for the same reason; Chapter~\ref{ch:data} explains it in full.

\subsection{Limitations of the Classical Approach}
\label{sec:classical_limitations}
The sliding-window method has three well-known weaknesses. All three feed into this thesis's move to a trained classifier, and the first two had to be handled directly in the pipeline.

\textbf{Weakness 1: the window length is set by hand, and the result depends on it.} Too short, and there are too few points in the window to measure the statistic reliably, so its trend is mostly noise. As a rough guide, the uncertainty in a lag-1 autocorrelation from $n$ points is about $1/\sqrt{n}$: for a 50-point window that is about $0.14$, close to the $0.2$--$0.4$ rise the method is trying to detect, so the signal is lost in the measurement error. This is a real problem here: after the onset ages were corrected, the pre-transition segments in the PANGAEA data range from about 100 to 650 points, and the shortest fall into this range (Chapter~\ref{ch:data}, Section~\ref{sec:label_correction} gives the numbers and the fix). Too long, and the window averages over the fast change just before the transition, and can stretch past the transition itself, breaking the assumption that everything inside it belongs to one stable regime.

\textbf{Weakness 2: detrending is unavoidable but can create a false signal.} Variance and autocorrelation are measured on the residuals left after the slow trend is removed, so the trend estimate affects the result. Smooth too strongly and real slow signal is removed with the trend, leaving residuals that are artificially anticorrelated. Smooth too weakly and some trend remains, raising the variance and autocorrelation on its own with no bifurcation involved. \citet{dablander2022} showed this for the deep-learning case: a classifier trained on residuals from one detrending filter partly learned to recognise the filter rather than the bifurcation, and misclassified when a different filter was used at test time. A preprocessing choice can quietly decide the result, for any method in this area.

\textbf{Weakness 3: the method detects change but not its kind.} This is the weakness that matters most here. Both a Fold bifurcation (a sudden, possibly permanent jump) and a Transcritical bifurcation (a smooth swap of stability) are preceded by the same rise in variance and autocorrelation. The two mean very different things, but on these statistics alone they cannot be told apart (Chapter~\ref{ch:theoretical_background}, Section~\ref{sec:fold_transcritical_challenge} explains why, and what extra information does separate them). A method built on variance and autocorrelation has nothing further to check; a classifier trained on many labelled examples can, in principle, learn the difference.

\section{The Deep Learning Shift: Bury et al. (2021)}
\label{sec:bury2021}
The method in this thesis is built most directly on \citet{bury2021}. That paper replaced the hand-picked statistic of the classical approach with a trained classifier, and showed it could be trained entirely on synthetic data and still work on real records. This section sets out what they did, what they found, and the three points where this thesis departs from them.

\subsection{The Core Idea}
The classical method commits in advance to one or two statistics -- variance, lag-1 autocorrelation -- and reads their trend. \citet{bury2021} instead treat early-warning detection as a classification problem: show a model many labelled examples of series approaching each kind of bifurcation, and many approaching nothing, and let it learn for itself what separates them.

The obstacle is data. Real tipping points are rare and mostly unlabelled, so a training set cannot be assembled from them. \citet{bury2021} get around this by generating the training set from simulated dynamical systems. Each training system is a pair of coupled equations,
\[
\dot{x} = \sum_{i} a_i\, p_i(x,y), \qquad \dot{y} = \sum_{i} b_i\, p_i(x,y),
\]
where the $p_i$ are every polynomial term in $x$ and $y$ up to third order, and the coefficients $a_i, b_i$ are drawn from a standard normal distribution, with half set to zero at random and the cubic terms forced negative to keep the solutions bounded. Each random system is simulated, checked for a stable equilibrium, and then passed to bifurcation-continuation software (AUTO-07P), which finds any fold, Hopf, or transcritical bifurcation along its equilibrium branch. Systems are generated until enough of each type has been collected, and the pre-bifurcation part of each simulation becomes one labelled training example.

This construction is what Chapter~\ref{ch:data}, Section~\ref{sec:sde_corpus} returns to, and it matters for how the results are read. The training systems are not the three textbook normal-form equations from Chapter~\ref{ch:theoretical_background}. They are a broad family of random systems that happen to pass through one of those bifurcations. The normal forms still describe each system near its bifurcation, by the Center Manifold Theorem, but the classifier learns from the full random systems, not the clean equations.

Two training sets were built this way: 500{,}000 series of length 500, and 200{,}000 series of length 1{,}500. Each residual series was detrended with a Lowess filter (span 0.2) and normalised by its own mean absolute value. The two sets were trained separately, giving a ``500-classifier'' for short records and a ``1{,}500-classifier'' for long ones. Four labels are used: fold, transcritical, Hopf, and a null class for systems with no approaching bifurcation.

\subsection{The Architecture and What It Scored}
The classifier is a CNN-LSTM: a convolutional layer reads short subsequences of the input and extracts local features, and an LSTM layer then reads that sequence of features in order and builds a summary of the whole series. \citet{bury2021} also tried a residual network, a convolutional network, and a recurrent network, and kept the CNN-LSTM because it scored highest on their training set.

On held-out synthetic data, an ensemble of ten 500-classifiers reached an F1 score of 84.2\% (precision 84.4\%, recall 84.2\%), and an ensemble of ten 1{,}500-classifiers reached 88.2\% (precision 88.3\%, recall 88.3\%). The longer input scored higher. This is the first evidence, from the paper this thesis builds on, that sequence length matters, and it is one reason this thesis runs every model at both lengths (Chapter~\ref{ch:data}, Section~\ref{sec:sde_corpus}).

They then compared the classifier against classical variance and lag-1 autocorrelation on eight test comparisons spanning six systems -- three model systems (a harvesting model, a predator-prey model, and a coupled disease-behaviour model) and three real datasets -- using the area under the ROC curve. The classifier beat both classical indicators on six of the eight. It predicted the correct bifurcation type in seven of the eight comparisons; the one exception was the thermoacoustic system, where the favoured-probability frequency for Hopf was only slightly higher than for fold. One behaviour is worth noting because it supports the whole approach: early in a series the classifier splits its probability roughly evenly across the three bifurcation types, and only past a certain point does it commit to one -- which is what a genuine early-warning signal emerging over time should look like.

\subsection{The Three Real-World Tests}
\citet{bury2021} applied the classifier to three real settings, none seen in training:

\begin{itemize}
  \item \textbf{Paleoclimate.} Seven of the eight abrupt climate shifts earlier analysed for early-warning signals by Dakos and colleagues.
  \item \textbf{A thermoacoustic experiment.} A laboratory Rijke tube driven to instability: 19 forced runs approaching the transition, and 10 held at a steady state.
  \item \textbf{Mediterranean anoxia.} 26 time series across eight anoxic events, from the three \citet{hennekam2020} cores -- the same records this thesis uses.
\end{itemize}

The classifier produced early-warning behaviour on all three despite being trained only on synthetic mathematics. That synthetic-to-real transfer is the strategy this thesis also follows (Chapter~\ref{ch:data}, Section~\ref{sec:dual_dataset}). The anoxia case is the direct overlap: this thesis runs its own evaluation on the same three cores, but on all sapropel events with corrected onset ages (Chapter~\ref{ch:data}, Section~\ref{sec:label_correction}), and across a much larger set of models, rather than on the subset used here.

\subsection{Where This Thesis Departs}
Three limits of the original study shaped this thesis.

\textbf{One architecture, one comparison.} \citet{bury2021} used a single design, CNN-LSTM, at two fixed input lengths (500 and 1{,}500), each applied to whichever real series matched its length. This thesis runs the same synthetic-to-real strategy across a much wider set -- 7 deep-learning models and 22 classical time series classifiers (Chapter~\ref{ch:methodology}) -- at both lengths, to see whether the CNN-LSTM's behaviour is a property of the task or of that one design. It largely carries over: across the 19 models that ran at both lengths, moving from 500 to 1{,}500 raises AUC by a median of about 3.6 points. The gain varies widely -- above +6 points for several models, and negative for one (MultiRocket, which does worse at the longer length). Chapter~\ref{ch:data} gives the per-model numbers.

\textbf{No explicit help for the Fold/Transcritical split.} As Section~\ref{sec:classical_limitations} set out, these two are the hard pair: both raise variance and autocorrelation, with no distinctive shape to separate them. \citet{bury2021} feed the classifier only the raw series and leave it to find the difference. This thesis adds the variance growth ratio channel (Chapter~\ref{ch:theoretical_background}, Section~\ref{sec:fold_transcritical_challenge}; Chapter~\ref{ch:data}, Section~\ref{sec:feature_engineering}) as a direct input for that one distinction. This thesis did not independently measure how well \citet{bury2021}'s own classifier separates fold from transcritical. Their main text reports only overall F1, precision, and recall figures; this thesis did not check their supplementary materials for a per-class breakdown, so no comparison on that point is claimed.

\textbf{Short real records must be padded.} A real sediment core is usually shorter than the classifier's fixed input length and has to be filled out to reach it. This thesis found the padding choice affects real-world accuracy, which is the subject of Section~\ref{sec:padding_gap}.

\section{Advancements in Time Series Classification}
\citet{bury2021}'s CNN-LSTM predates a decade of progress in general-purpose time series classification (TSC), summarised in \citet{bagnall2017bakeoff}'s large-scale comparison: 18 published classifiers plus 2 baseline classifiers (1-NN DTW and Rotation Forest), run across 85 datasets from the UCR archive, with only 9 significantly beating both baselines. The paper organises classifiers into six families: distance-based, difference-based, dictionary-based, shapelet-based, interval-based, and ensemble. The clear overall winner was COTE, an ensemble across several of these representations, on average about 8\% more accurate than either baseline; the next two best individual classifiers were a shapelet method (Shapelet Transform, second overall) and a dictionary method (BOSS, third overall), and three interval-based classifiers (TSF, TSBF, LPS) all beat both baselines with no significant difference between them. Kernel-based methods -- the ROCKET family this thesis relies on heavily -- did not exist yet in 2017 and so were not part of this comparison; this thesis draws on them as the next major development along the same line. The bake-off's task is also different from this thesis's own: general-purpose sequence classification (gesture recognition, sensor data, and similar), not bifurcation-type classification. What carries over is not a performance number but the practice of benchmarking broadly across model families instead of committing to one architecture, and the specific families -- dictionary-based, shapelet-based, interval-based, and, from later work, kernel-based -- that this thesis's own roster (Chapter~\ref{ch:methodology}) is built around.

\subsection{Why Not Just Recurrence}
An LSTM reads a series one step at a time and keeps a single hidden-state vector meant to summarise everything seen so far. Two things make this awkward for a long, noisy paleoclimate record. First, that vector is a fixed-size bottleneck: whatever mattered many steps ago has to survive being repeatedly overwritten, so a slow signal spread across a 1,500-point series competes for space with a lot of step-to-step noise. Second, training an LSTM over a long series means propagating a gradient back through as many steps, and that gradient tends to shrink or grow at each one -- the vanishing/exploding gradient problem -- which makes long-range dependencies harder to learn than short-range ones. A convolutional or kernel-based model instead looks at many fixed-size windows of the series at once and never has to carry information step by step across the whole length, which sidesteps both problems. This is part of the motivation for testing kernel-based and interval-based families here rather than relying on recurrence alone.

\subsection{Kernel-Based Methods: ROCKET and its Successors}
\citet{dempster2020} introduced ROCKET (RandOm Convolutional KErnel Transform): generate a large number of random convolutional kernels (10{,}000 by default), each with a randomly drawn length, weight, bias, and dilation, convolve the series with every one, and reduce each kernel's output to two numbers -- the maximum value and the proportion of positive values (PPV) -- then feed the resulting 20{,}000 features into a fast linear classifier. No kernel is ever trained by gradient descent; using thousands of random kernels is enough on its own.

\citet{dempster2021minirocket} (MiniRocket) removes almost all of that randomness. It uses a small, fixed set of 84 length-9 kernels, each built from just two integer weights (three $+2$s and six $-1$s, chosen so the convolution needs only addition), and keeps only the PPV feature, dropping ROCKET's maximum. The kernels' biases come from quantiles of the training data itself rather than being drawn at random, which is what makes the method ``almost,'' rather than fully, deterministic. The result runs over an order of magnitude faster than ROCKET with essentially the same accuracy.

\citet{tan2022multirocket} (MultiRocket) keeps MiniRocket's fixed kernels but adds three further pooling statistics alongside PPV -- the mean of the positive values, the mean of their positions in the window, and the length of the longest run of positive values -- and applies the whole transform twice, once to the raw series and once to its first-order difference. This raises the feature count from ROCKET's 20{,}000 to 50{,}000, aimed at capturing more about each kernel's response than just how often it activates.

All three are part of the model roster used here. Whether the kernels are drawn at random (ROCKET) or fixed but varied (MiniRocket, MultiRocket), the spread of kernel dilations gives this family sensitivity to patterns at many different time scales in one pass, which suits a signal like the CSD ramp that can show up over very different fractions of a record depending on how close to the tipping point a segment starts.

\subsection{Dictionary and Interval Methods}
WEASEL 2.0 \citep{schafer2023weasel2} converts sliding windows of a series into discrete "words" using the Symbolic Fourier Approximation, then classifies on which words occur and how often -- a bag-of-patterns approach built in the frequency domain. Its departure from earlier dictionary methods, including BOSS, is to draw the window length and the dilation between sampled points at random for each word rather than fixing them in advance, so the dictionary covers several time scales without the user choosing them.

Time Series Forest \citep{deng2013tsf} takes a different route: each tree in the ensemble is given $\sqrt{m}$ random intervals of the series (for a series of length $m$), computes three summary statistics on each -- mean, standard deviation, and slope -- and is trained on the resulting $3\sqrt{m}$ features; the forest classifies by majority vote. \citet{bagnall2017bakeoff} found TSF, alongside two related interval methods, significantly more accurate than both their benchmark classifiers, with no significant difference between the three, and recommended TSF specifically for its simplicity.

Both are part of the model roster used in this thesis, alongside several further interval, shapelet, and dictionary-based methods described in Chapter~\ref{ch:methodology}.

\subsection{Closely Related Recent Studies}
Two papers apply almost exactly this thesis's method -- deep learning trained on synthetic bifurcation data, tested on real anoxia records -- to the same underlying real dataset used here.

\citet{babazadeh2025} train a CNN-LSTM classifier on synthetic fold/transcritical/Hopf/null data built directly from the normal forms (rather than Bury et al.'s wider random-system library) and test it, among other things, on the same Mediterranean anoxia cores from \citet{hennekam2020} used in this thesis. Two of their findings are directly relevant here. First, their classifier's predictions on the real anoxia data favoured a Hopf-type bifurcation, in tension with the fold-type onset expected on domain grounds. Second, they tested robustness to autocorrelated ("coloured") training noise by replacing white noise with red noise of varying strength, and found validation accuracy stayed within a narrow 83--85\% band across the whole range tested, with AUC scores of 0.9 to 1 on coloured test series for classifiers trained on either kind of noise.

\citet{ma2025sdml} -- co-authored by \citet{bury2021}'s own lead author -- take a different route to the same real cores: rather than simulating dynamical systems from first principles, their surrogate data-based machine learning (SDML) approach generates training data by applying surrogate methods (AAFT/IAAFT) directly to real historical time series, then trains several architectures (SVM, LSTM, CNN, Multi-Head CNN) and reports, per real-world test system, whichever architecture scored best on its synthetic validation set: SVM for the MS66 core (F1 = 1.0), LSTM for 64PE406E1 (F1 = 0.99), and CNN for MS21 (F1 = 0.99). The LSTM architecture they describe is the same one this thesis reproduces as its own \texttt{lstm} model (Chapter~\ref{ch:methodology}, Section~\ref{sec:dl_architectures}), confirmed by comparing their stated layer sequence against this thesis's own \texttt{models/lstm.py} directly. These reported F1 scores are a synthetic-validation metric, the same kind of ground-truth-known figure as this thesis's own Table~\ref{tab:zenodo_results}-style results, not a real-data accuracy figure -- their paper states explicitly that the real anoxia data used to generate their ROC curves is not involved in producing this F1 number, so it is not compared here against any of this thesis's own real-data results.

None of these three papers' real-data findings are re-derived or re-confirmed independently by this thesis -- they are noted here because this thesis's own empirical analysis on the same real cores reaches a similar qualitative conclusion about Hopf-type predictions dominating, from a much larger set of independently-trained models (Chapter~\ref{ch:results}), which is worth reading against these prior results rather than in isolation.

\section{Left-Censored Data and Padding}
\label{sec:padding_gap}
Every model here needs an input of a fixed length: 500 or 1{,}500 points, depending on which one it was trained at. A real sediment core almost never has that many usable points before the transition. After the sapropel onset ages were corrected (Chapter~\ref{ch:data}, Section~\ref{sec:label_correction}), the seven segments this thesis uses run from about 100 to 650 points -- all but one shorter than even the smaller model size, and every one shorter than the larger model size of 1{,}500 (Chapter~\ref{ch:data}, Figure~\ref{fig:segment_lengths}). This missing length is not random: the record simply does not go back far enough, at the resolution needed, to fill it. Something has to be added to reach the required length.

\subsection{Padding Choice Matters, More Than It Might Seem}
The simplest fix is zero-padding: add zeros to the front of the record until it reaches the target length. \citet{bury2021} do the same in their own training: they pad each training example with zeros, either on both ends or on the left only, before feeding it to the model. This thesis uses zero-padding throughout: every result reported here comes from \texttt{pad\_mode: "zero"} in \texttt{config.yaml}.

Zero-padding creates a sharp jump: the input goes straight from exact zero to the real signal. The model cannot tell that this jump is just an artefact of preprocessing, not something that actually happened in the data. This is a general reason for caution, not a specific effect this thesis has measured. Two other ways to fill the gap avoid the sharp jump. One repeats the first real value instead of using zeros, so the level does not jump. The other reflects the real signal back across the boundary, so the level and the direction it was moving both carry over smoothly.

There is a real precedent for a preprocessing choice like this mattering, even though it was for a different choice. \citet{dablander2022} showed that the \citet{bury2021} classifier is sensitive to which detrending filter it is given. A classifier trained on residuals from a Lowess filter seems to have partly learned to recognise that filter, not just the bifurcation, and they recommend using the same filter at test time that was used in training. This thesis did not find a discussion of padding in their paper; the idea that padding could matter in the same way is this thesis's own extension of their finding, not a claim they made. This thesis's pipeline implements the two alternatives -- repeating the first value, and reflecting the signal -- alongside zero-padding, in \texttt{src/data\_common.py}.

It is worth being precise about what was and was not tested here. Every result in this thesis used zero-padding. The other two padding modes exist in the code but were never used to produce a reported result. What \emph{was} tested is a different question: not which padding value to use, but how much padding a training example gets while training, and whether that amount should be adjusted to match what real short segments need at test time. Chapter~\ref{ch:data}, Section~\ref{sec:preprocessing} covers this directly: the amount of padding used during training was capped and randomised, and both limits were tightened so that training matches what the real cores actually need. Testing zero-padding against the other two padding modes on the real cores is a natural next step, now that both are implemented -- but it has not been done yet, so it is future work, not a finished contribution.

\section{Gap in the Literature and This Thesis's Contribution}
Three papers trace one line of work. \citet{scheffer2009} explained why critical slowing down should be a generic warning sign. \citet{dakos2012} turned that idea into a fixed, repeatable statistical method, but one that depends on choices like window length and detrending bandwidth. \citet{bury2021} replaced the hand-picked statistic with a trained classifier, which removes those choices, but tested only one architecture, and each of its three real-world tests was small: a handful of climate records, one lab experiment, and a subset of anoxia events. This thesis's contribution sits at the next step: many architectures instead of one, and a larger, corrected version of the same real anoxia data Bury et al. used. The closely related \citet{babazadeh2025} sits alongside this thesis at a similar step, from a different direction; Section~\ref{sec:padding_gap} already discussed one difference between the two (padding), and Table~\ref{tab:related_work_comparison} compares all three directly.

\begin{table}[htbp]
\centering
\caption{The three deep-learning approaches to this problem, compared directly.}
\label{tab:related_work_comparison}
\begin{tabular}{p{2.4cm}p{3.4cm}p{3.4cm}p{3.6cm}}
\toprule
 & \textbf{Bury et al. (2021)} & \textbf{Babazadeh et al. (2025)} & \textbf{This thesis} \\
\midrule
Architecture(s) & CNN-LSTM & CNN-LSTM & 7 DL + 22 classical TSC \\
Training data & Random-system library, verified by continuation software & Normal forms simulated directly & Bury's released random-system library (Chapter~\ref{ch:data}) \\
Sequence lengths & 500 and 1500 & Training series length 100 & 500 and 1500, across every model \\
Noise robustness & White noise only & White vs.\ coloured noise, tested & Not tested (flagged as a limitation) \\
Real test data & Paleoclimate transitions, thermoacoustic rig, Mediterranean anoxia (subset) & Thermoacoustic rig, Mediterranean anoxia & Mediterranean anoxia only, full 3-core set, corrected onset ages \\
Type-diagnosis reliability across independent models & Not measured & Not measured & Measured directly (Chapter~\ref{ch:results}) \\
\bottomrule
\end{tabular}
\end{table}

\subsection{The Contributions of This Thesis}
\begin{enumerate}
    \item \textbf{A much broader model comparison.} Rather than one architecture, this thesis benchmarks 7 deep-learning models and 22 classical time series classifiers (Chapter~\ref{ch:methodology}) on the same synthetic-training, real-testing protocol, at both sequence lengths. As far as this thesis is aware, no earlier study has run a comparison at this scale on this specific task.
    \item \textbf{An explicit feature for the Fold/Transcritical distinction.} The variance growth ratio channel (Chapter~\ref{ch:theoretical_background}, Section~\ref{sec:fold_transcritical_challenge}) gives the classical (non-deep-learning) models in the roster a channel tracking how far the windowed variance has grown relative to its baseline value -- a log-scaled ratio, not just the variance level itself. The deep-learning models receive only the raw residual and have to learn any such distinction from that alone. Because a Fold's variance accelerates while a Transcritical's grows steadily, the two leave differently-shaped traces in this channel, which is the one distinction classical EWS structurally cannot make.
\end{enumerate}
